% 17-agents.tex: программные агенты, доказательство личности человека и будущий интернет.
\section{Программные агенты, доказательство того, что ты человек, и будущий интернет}
\label{sec:agents}

Этот раздел наносит на карту, что устанавливают механизмы стороны держателя, когда предъявляющий
удостоверение может оказаться агентом, а не человеком: четыре механизма, что каждый доказывает, чему
каждый отказывает и какой остаток не достигает ни один из них.

\subsection{Задача, которая приближается}

Когда программные агенты способны порождать неограниченное число учётных записей, материалов,
операций, сообщений и мнимых участников, дешёвый признак того, что запрос пришёл от человека,
перестаёт доставаться даром. Многим цифровым системам придётся различать, не узнавая лишнего,
одного подотчётного человека; одного уполномоченного агента, действующего за этого человека или
за учреждение; безымянную машинную активность; сдвоенные или синтетические личности; и агентов
организаций, действующих под полномочиями учреждения. Всякий простодушный способ получить этот
признак выдаёт личность: учётные записи, номера телефонов, платёжные средства, биометрия. Вопрос о
личности меняет форму. Это уже не только вопрос \emph{кто ты}, но и \emph{что ты за деятель, чьи
полномочия ты осуществляешь, что тебе дозволено делать и как мне всё это проверить, ничего больше
не узнав}. Это и есть обращённый в будущее вопрос, который настоящий документ считает самым
важным из тех, на которые можно было бы научить отвечать ядро Поляриса.

\subsection{Что уже существует и указывает в эту сторону}

Работают шесть составляющих доказательства того, что ты человек, не раскрывающего личности.

\begin{itemize}
  \item \textbf{Единственность при выдаче.} С3, один действующий токен на человека, есть частичный
    уникальный индекс в базе данных и модель, проверяемая при каждой отправке. Устанавливает она уникальность удостоверения в пределах записанной личности, а не
    уникальность человека: у одной строки \entity{ФизическоеЛицо} не больше одного действующего удостоверения,
    и никакой код приложения не способен это обойти. Ничто ни в базе данных, ни над ней не замечает
    одного человека, зачисленного двумя строками; это решается у стола зачисления суждением оператора,
    записанным как свидетельство (раздел~\ref{sec:privacy}).
    Поэтому С3 не есть устойчивость к атаке Сивиллы. \sImpl
  \item \textbf{Принадлежность без личности.} Доказательство принадлежности эпохе доказывает, что
    действительный токен существует в зафиксированном множестве органа, при открытых входах из
    одних лишь корня, эпохи, контекста, одноразового числа и нуллификатора с областью; проверка с
    нулевым разглашением не записывает никакого токена вовсе; множество анонимности есть целая
    эпоха, и, как показали измерения, ничто не задаёт ему нижнего предела выше единицы. \sImpl
  \item \textbf{Отзыв без называния.} Отозванный токен выпадает из корня следующей эпохи, так что
    доказательство принадлежности истекает вместе с эпохой, и ни одна доверяющая сторона не
    узнаёт, кто был отозван. \sImpl
  \item \textbf{Ключ, которым распоряжается держатель.} Удостоверение может нести открытый ключ
    держателя, который эмитент к нему привязывает, так что владение файлом перестаёт быть владением
    удостоверением, а верификатор вправе потребовать подпись самого человека, а не копию его
    удостоверения. \sImpl
  \item \textbf{Доказывание на устройстве держателя.} Множество листьев эпохи публикуется,
    подписанным и ограниченным; каждый запрашивающий получает одинаковые байты; держатель находит
    свой лист и строит путь у себя. Эмитенту никогда не сообщают, какой именно участник
    доказывает. \sImpl
  \item \textbf{Нуллификатор для каждого верификатора.} Доказательство несёт
    \code{Посейдон(секрет\,||\,область\,||\,эпоха)}, постоянный для одного человека в пределах
    области и эпохи одной доверяющей стороны и не соотносимый между областями. \sImpl
\end{itemize}

\subsection{Один человек однажды: что работает и чего это стоит}

Доказательство принадлежности не даёт переиграть одно доказательство, но само по себе не мешает
одному и тому же человеку изготовить тысячу свежих доказательств для одного сайта и завести
тысячу учётных записей. Правилу \emph{один человек, однажды, здесь} нужен нуллификатор с
областью: ещё один открытый вход, хеш секрета держателя, области доверяющей стороны и эпохи, так
что сайт может отвергнуть второе доказательство от того же человека, а два сайта не могут связать
свои нуллификаторы. Это открытый вход схемы, а утверждение, что две доверяющие стороны, объединив
всё, чем владеют, всё равно не свяжут держателя, есть формальная модель, проверяемая вместе с
конфигурацией, на которой она обязана упасть. \sImpl

Лист раскрывается внутри схемы, и причина структурная. Нуллификатор рядом с листом, который схема не
может раскрыть, доказывает лишь то, что доказывающий знает какое-то число: ничто не привязывает их к
одному секрету, так что доказывающий мог бы соединить лист любого участника с нуллификатором по
собственному выбору, и каждое перечисленное выше свойство было бы утверждением, а не доказательством.
Привязать оба к одному секрету значит раскрыть лист в схеме, где прообраз \code{ША3-256} стоит тысячи
ограничений, а Посейдон около сотни. Поэтому лист есть \code{Посейдон(секрет\,||\,контекст)},
обязательство, которое схема раскрывает, а секрет держателя выводится через \code{ША3-256} вне её.

Два условия, которые ставит призвание, были перенесены в реализацию, а не оставлены намерениями.
Область нуллификатора задаётся каждой доверяющей стороной и каждой эпохой и никогда не бывает
глобальной, потому что глобальный нуллификатор есть постоянный идентификатор, а значит, ровно тот
исход слежки, который отвергает конституция; и он меняется с каждой эпохой, так что
принадлежность никогда не становится постоянным псевдонимом. Учёт доверяющей стороны привязан к её
собственной области и к этой эпохе и не должен их переживать.

Две границы остаются и изложены, а не закрыты. Нуллификатор не прячет держателя от
\emph{эмитента}, который выводит каждый лист, чтобы построить дерево, и потому способен вычислить
нуллификатор любого участника в любой области; свойство действует между доверяющими сторонами. А
слой предъявления по-прежнему предъявляет полное удостоверение на всяком пути, кроме пути с нулевым
разглашением: верификатор, глядящий на обыкновенное удостоверение, видит устойчивое значение
токена, подпись эмитента и ключ держателя, а верификатор на пути совместимости видит девять
подписанных эмитентом значений, одинаковых повсюду. Попарный вывод ограничивает то, что верификатор
\emph{хранит}, то есть тот материал, который объединяют, продают, истребуют по повестке и теряют
при утечках; путь с нулевым разглашением ограничивает то, что ему \emph{показывают}, и лаборатория
измерила эту границу и её слепые пятна (раздел~\ref{sec:privacy}). Верификатор сообщает, что из
двух дало данное предъявление, и со времени первой находки лаборатории сперва проверяет посылку и
лишь потом это говорит. Технологии ослеплённых и одноразовых предъявлений, которые устранили бы
остаток, отложены, а не запланированы, в разделе~\ref{sec:research}; ни одна не реализована. \sFut

\subsection{Передача полномочий агенту}

Это отношение есть цепочка, и она работает. Человек владеет основным удостоверением и ключом,
который эмитент к нему привязал. Этим ключом человек подписывает полномочие, дающее агенту узко
ограниченную власть: названные действия, изложенные пределы, срок и дескриптор отзыва,
принадлежащий одному этому полномочию. Агент подписывает каждое совершаемое действие, называя
действие и собственное одноразовое число службы. Служба проверяет автономно, что полномочие
подлинно и не истекло, что оно покрывает именно это действие, что оно не отозвано и что оно
сцеплено с ключом держателя, который эмитент привязал к настоящему удостоверению. \sImpl

\figfile{delegation}{Цепочка передачи полномочий в том виде, как она построена. Человек держит
основное удостоверение и ключ, привязанный к нему эмитентом. Этим ключом он подписывает полномочие,
называющее действия, пределы, срок и дескриптор отзыва; агент подписывает каждое действие против
собственного одноразового числа службы; служба проверяет цепочку автономно. Единственное звено,
нарисованное пунктиром, есть остаток: при обыкновенной привязке держателя служба видит ещё и
значение токена удостоверения, так что соотносимость, которая ей остаётся, ограничена, а не
отсутствует.}{fig:delegation}

То, ради чего задумано это устройство, видно по тому, в чём оно отказывает. Передать агенту само
удостоверение значит дать ему всё, что может сделать человек, навсегда, с возможностью отзыва
лишь через отзыв самого человека. Полномочие во всём противоположно этому: \code{действия} и
\code{пределы} лежат \emph{внутри} подписанного утверждения, так что полномочие, расширенное по
дороге, падает, а не проходит незаметно; пустой перечень действий не разрешает ничего, и способа
написать \emph{всё} нет; ключ предела, которого служба не знает, отвергается, а не пропускается
мимо, поскольку полномочие, говорящее \code{макс\_переводов: 3} службе, которая об этом пределе
никогда не слышала, не должно читаться как неограниченное. И агент обязан подписывать: без этого
звена полномочие есть токен на предъявителя, и всякий, кто скопирует его по дороге, становится
агентом. Два из этих отказов, исчерпанный счётчик использований и сумма сверх потолка, суть та
граница, которая делает передачу полномочий чем-то иным, нежели токен на предъявителя; контракт
соответствия включает случаи с плохой подписью полномочия или его отзыва, а мутационное учение
показывает, что тесты заметили бы утрату любой из двух границ. \sImpl

Отзыв есть то свойство, которое стоит назвать дважды. Отзыв подписывается тем же ключом, которым
подписано полномочие, так что опубликовать байты может кто угодно, а прекратить полномочие может
лишь держатель. К эмитенту никто не обращается, он никогда не узнаёт, что полномочие существовало,
и удостоверение человека всё это время остаётся нетронутым. Человек может прекратить власть своего
агента, не спрашивая разрешения у органа, выдавшего его личность, и так, что этот орган этого не
увидит. Отзыв несёт идентификатор полномочия и момент, и больше ничего. Причины он нести не будет:
поле, записывающее, \emph{почему} полномочие прекращено, есть поле, которое принуждающий может
потребовать заполнить или оставить пустым, и в обоих случаях оно превращает отзыв в сигнал о
человеке.

То, что узнаёт служба, ограничено, и граница сообщается, а не предполагается. При обыкновенной
привязке держателя служба видит ещё и значение токена удостоверения, поэтому вердикт называет
соотносимость \code{открыто}; лишь цепочка, переданная по пути с нулевым разглашением, чей
дескриптор есть нуллификатор с областью, а протокол не показывает ничего одинакового у разных
верификаторов, называется \code{ограничено}. Полномочие прячет человека от \emph{действий} агента.
Оно не прячет удостоверение от службы.

\subsection{Будущий интернет, конкретно}

Слово Веб3 употребляется здесь в одном узком смысле и не имеет никакого отношения к валюте: это
интернет, в котором людям, самостоятельным агентам, учреждениям, кошелькам, распределённым службам
и независимо эксплуатируемым сетям нужно переносимое криптографическое доверие, не проходящее
через одну компанию или одно государство. Против такой среды следующая таблица отмечает, что
Полярис мог бы предоставить и каково сегодня состояние каждого пункта.

{\footnotesize
\setlength{\tabcolsep}{4pt}\renewcommand{\arraystretch}{1.12}
\rowcolors{2}{tabstripe}{white}
\begin{xltabular}{\textwidth}{>{\RaggedRight\arraybackslash\hspace{0pt}}p{1.85in} >{\RaggedRight\arraybackslash\hspace{0pt}}X >{\RaggedRight\arraybackslash\hspace{0pt}}p{1.8in}}
\caption{Чего будущая среда потребовала бы от Поляриса и каково состояние каждого пункта.}\label{tab:web3}\\
\toprule
\rowcolor{tabheader}\thead{Среда просит} & \thead{Что есть у Поляриса} & \thead{Состояние} \\
\midrule
\endfirsthead
\toprule
\rowcolor{tabheader}\thead{Среда просит} & \thead{Что есть у Поляриса} & \thead{Состояние} \\
\midrule
\endhead
\midrule\multicolumn{3}{r}{\itshape продолжение на следующей странице}\\
\endfoot
\bottomrule
\endlastfoot
Переносимые удостоверения & Подписанный файл, который проверяет любая стандартная библиотека; четыре пакета в публичных реестрах & \sImpl \\
Независимо проверяемые доказательства & Отделяемый верификатор, два КСР, набор проверок соответствия и верификатор ОИД4ВП, которому предъявлял сторонний кошелёк & \sImpl\newline\sWit\newline\sExt \\
Федерация учреждений & Корни у каждого органа, направленные аттестации, опубликованные ленты & \sImpl\newline\sExt \\
Обнаружение служб & Подписанный реестр & \sImpl \\
Подписи документов с долгосрочной действительностью & Подписанные контейнеры со свидетельствами, закреплёнными в момент подписания & \sImpl \\
Прозрачное поведение власти & Три журнала, свидетели, доказательства двоесказательства, публичный отчёт о чтениях по ордеру & \sImpl\newline\sExt \\
Автономная проверка & Утверждения о состоянии, подшитые к предъявлениям; карта, подписывающая вызов устройства & \sImpl \\
Постквантовое доверие & МЛ-ДСА при проверяемом пути миграции, хеш-основанные доказательства, гибридная кромка; классические остатки названы; во что обошёлся бы взлом, записано & \sImpl \\
Отношения полномочий между человеком, учётной записью и агентом & Входной токен, субъект которого свой для каждой доверяющей стороны, и полномочие агента & \sImpl \\
Единственность без раскрытия личности & Доказательства принадлежности с нуллификатором для каждого верификатора и каждой эпохи; измерено против сговорившейся пары & \sImpl \\
Передача агентам полномочий с ограниченной областью & Полномочие, подписанное держателем: действия, пределы, срок, дескриптор отзыва; его границы проверены мутацией & \sImpl \\
Отзыв агента без разрушения личности человека & Отзыв, подписанный держателем; к эмитенту не обращаются, удостоверение не тронуто & \sImpl \\
Предъявление, не показывающее верификатору ничего устойчивого & Только путь с нулевым разглашением, в пределах множества анонимности; обыкновенное предъявление и предъявление совместимости показывают устойчивые значения, подписанные эмитентом & \sImpl\newline\sFut \\
Разговор на форматах других кошельков & Проверенный ОИД4ВП с СД-ЖВТ ПУ; отображение в ИСО 18013-5, подпись эмитента в котором не проверит ни один стандартный считыватель; проверяемое удостоверение о результате проверки & \sImpl\newline\sWit \\
\end{xltabular}
}
